\section{Formal benchmarks for the mental-models treatment}
\label{sec:annex-p19-rationality-2-3}

These finite derivations justify declared semantic and numerical examples.
They do not establish that participants use the corresponding internal
representations, or quantify the effectiveness of a training intervention.

\begin{annexitem}{annex:p19.r23.nonvacuity}{thm:p19-r23-nonvacuity}{Ordinary and nonvacuous consequence}
\begin{proof}
Both definitions require that every valuation satisfying all the premises
satisfy the conclusion. When \(\mathcal M(\Gamma)\) is nonempty,
the extra nonemptiness requirement is already satisfied, so the definitions
agree. When it is empty, the universal assertion about its members is
true for every conclusion, but the nonemptiness conjunct is false for
every conclusion. The existential requirement for a live possibility
also fails. No Boolean valuation makes \(c\) and \(\neg c\) both
true, establishing the example independently of the value of \(t\).
\end{proof}
\end{annexitem}

\begin{annexitem}{annex:p19.r23.possibility}{thm:p19-r23-possibility}{The two kinds of possibility judgment}
\begin{proof}
The complete valuations satisfying \(c\lor t\), in coordinate order
\((c,t)\), are \((1,0),(0,1),(1,1)\). The first is a witness to
\(c\), the second to \(t\), and the third to \(c\land t\).
Adding \(\neg c\) leaves only \((0,1)\), so there is no witness
to \(c\) or \(c\land t\) after that addition.
The exactly-one valuations \((1,0),(0,1)\) both satisfy the inclusive
disjunction. Conversely \((1,1)\) satisfies the inclusive assertion
but not the exactly-one assertion, proving the stated Boolean contrast.

For the modal statements, take one world \(w\) with accessibility
relation \(\{(w,w)\}\), make \(c\) false and \(t\) true at
\(w\). The premise \(c\lor t\) is true there. Its only accessible
world makes \(c\) false, so both \(\Diamond c\) and
\(\Diamond(c\land t)\) are false. This is a reflexive, transitive
and symmetric frame, hence in particular a normal K, T, S4 and S5
countermodel under their usual Kripke semantics. The existence of some
valuation compatible with a premise family does not impose that an
arbitrary actual world's accessibility relation contains that valuation.
The finite live-possibility operation can lose a former witness when
premises are added; ordinary modal entailment with a fixed interpretation
has its different, monotone premise relation.
\end{proof}
\end{annexitem}

\begin{annexitem}{annex:p19.r23.xor}{thm:p19-r23-xor}{Complete valuation check of the illusion}
\begin{proof}
Exclusive disjunction is true when its two arguments differ. For
\((c,t)=(0,0),(0,1),(1,0),(1,1)\), respectively,
\(c\mathbin{\oplus}t\) has values \(0,1,1,0\). Because
\(\neg t\) reverses the second argument,
\(c\mathbin{\oplus}\neg t\) has values \(1,0,0,1\).
No row makes both assertions true. The premise-compatible set is empty;
the nonvacuous criterion in \cref{def:p19-r23-live-models} therefore
licenses neither specified conclusion. Under ordinary classical
consequence both would follow vacuously, which is why the criterion
must be named rather than left implicit.
\end{proof}
\end{annexitem}

\begin{annexitem}{annex:p19.r23.joint}{thm:p19-r23-joint}{Cell reconstruction and all three admissibility tests}
\begin{proof}
The four intersections partition the sample space. Additivity gives
\(P(A\neg B)=a-q\) and \(P(\neg A B)=b-q\), leaving
\(P(\neg A\neg B)=1-a-b+q\). Their displayed sum is one
for all real \(a,b,q\), independently of signs. With \(a,b\in[0,1]\),
nonnegativity of the four cells is equivalent to
\(q\geq0,q\leq a,q\leq b,q\geq a+b-1\), precisely the
stated interval. If these inequalities hold, defining a four-point
probability space with these cells gives the required events and
marginals, proving sufficiency as well as necessity. Each cell is at most
one because it is nonnegative and the total is one.

By inclusion-exclusion, \(u=a+b-q\). Substitution in the upper bound
on \(q\) gives \(u\geq\max(a,b)\), while substitution in its
lower bound gives \(u\leq\min(1,a+b)\); the construction just used
also proves sufficiency. For \(a>0\), the ordinary conditional ratio
gives \(r=q/a\), hence \(q=ar\), so exactly the same four
nonnegativity inequalities apply. If \(a=0\), any admissible joint
distribution has \(q=0\), and the ratio \(0/0\) does not define
\(P(B\mid A)\). No arbitrary value for that conditional has been
inferred from the marginals. If \(a,b>0\), divide the same intersection
\(q=ar\) by \(b\) to obtain \(r'=ar/b\). If \(b=0\),
the converse conditional's ratio is likewise undefined. For unequal
\(a,b\), the arithmetic mean lies strictly between them. It therefore
exceeds \(\min(a,b)\), violating the intersection upper bound, and
is below \(\max(a,b)\), violating the union lower bound. The worked
negative-cell triple gives
\(.5+.3-.5+.4-.5+1-.3-.4+.5=1\); its two negative cells
nevertheless exclude a probability interpretation.
\end{proof}
\end{annexitem}

\begin{annexitem}{annex:p19.r23.kernel}{thm:p19-r23-kernel}{Admissible conditional branches and their converse}
\begin{proof}
Every factor in each displayed cell lies in \([0,1]\), so the cells
are nonnegative. Summing the first two gives \(a\), and summing the
last two gives \(1-a\), so the total is one. For \(a>0\),
the first cell divided by the first-row sum is \(r\); for \(a<1\),
the third cell divided by the second-row sum is \(s\).

Conversely take any four nonnegative cells
\((q_{11},q_{10},q_{01},q_{00})\) summing to one. Set
\(a=q_{11}+q_{10}\). If \(0<a<1\), set
\(r=q_{11}/a\) and \(s=q_{01}/(1-a)\); these numbers lie in
\([0,1]\) and substitution recovers every cell. If \(a=0\),
nonnegativity forces \(q_{11}=q_{10}=0\). Set \(s=q_{01}\)
and choose any \(r\in[0,1]\). The formula then recovers the
original cells but does not identify \(r\). If \(a=1\),
\(q_{01}=q_{00}=0\); set \(r=q_{11}\) and choose any
\(s\in[0,1]\). The same conclusion holds for this unused branch.
These cases exhaust all possibilities and respect the undefined ordinary
conditional at a null conditioning event.
\end{proof}
\end{annexitem}

\begin{annexitem}{annex:p19.r23.chances}{thm:p19-r23-chances}{The two finite counting examples}
\begin{proof}
The disease counts sum to ten. Exactly five chances have the symptom,
and exactly three have both symptom and disease. Equal chance weight
therefore gives \(P(S)=5/10>0\) and
\(P(D\mid S)=(3/10)/(5/10)=3/5\).

For the marble assertion, when \(R=1\), exclusivity requires
\(G\land B=0\), giving \(100,110,101\). When \(R=0\),
it requires \(G=B=1\), giving \(011\). These cases include
every satisfying triple and no others. Among them only \(101\)
satisfies \(R\land B\). It is therefore a witness to possibility
and receives one of four equal weights under uniform counting.
Alternatively give weight one to \(100\) and zero to the other
three triples. This is a probability distribution on the same
possibility set and gives the target probability zero. A stipulation
of positive mass for every satisfying triple would exclude this
second measure; possibility alone does not make that stipulation.
\end{proof}
\end{annexitem}

\begin{annexitem}{annex:p19.r23.cards}{thm:p19-r23-cards}{Sufficiency and necessity for complete card checking}
\begin{proof}
A counterexample must combine the letter E with a number other than 2.
The E card can be a counterexample if its hidden number is not 2.
The 3 card can be a counterexample if its hidden letter is E.
Turning both reveals whether either
counterexample is present. The F card cannot be a counterexample because
its known letter is not E, and the 2 card cannot be one because
its known number is 2. These cases exhaust the four cards, proving
sufficiency.

If E is not turned, two arrangements can agree on every inspected
face while differing only in whether its hidden face is 2 or 3. Choose
all other hidden faces to obey the rule. The rule then holds in one
arrangement and fails in the other, so uninspected E prevents a
conclusive decision. If 3 is not turned, make its hidden face F in one
arrangement and E in the other, again fixing all other cards to obey
the rule. The same indistinguishability argument applies. Thus each of
E and 3 is necessary when no other information restricts the hidden
faces. This proof makes no claim about a different objective that uses
prior probabilities, unequal costs or nonconclusive evidence gain.
\end{proof}
\end{annexitem}

\begin{annexitem}{annex:p19.r23.reversal}{thm:p19-r23-reversal}{Train invariant, termination and transfer count}
\begin{proof}
For \(n=0\), the specified no-operation branch has the required
empty output. Let \(n\geq1\). If \(n>1\), the initial bulk
transfer is legal by the terminal-block rule and gives
\(a_1[a_2\cdots a_n]\epsilon\). If \(n=1\), this state is
already the initial state without a bulk transfer. The first one-car
right transfer gives \(\epsilon[a_2\cdots a_n]a_1\), establishing
the invariant at \(k=1\).

Suppose the invariant holds at \(k<n\). The siding is nonempty
with first car \(a_{k+1}\). Moving that initial one-car block to
the empty left word gives
\(a_{k+1}[a_{k+2}\cdots a_n]a_k\cdots a_1\).
Moving the left word's terminal one-car block to the beginning of
the right word gives
\(\epsilon[a_{k+2}\cdots a_n]a_{k+1}a_k\cdots a_1\),
the invariant at \(k+1\). Each operation is permitted and preserves
car order within its moved block. The argument does not require
distinct printed labels.

The number of siding cars is \(n-k\), a nonnegative integer that
decreases by one each iteration. Thus exactly \(n-1\) iterations
occur, and the state at \(k=n\) is the asserted final state.
For \(n=1\) there is only the first right transfer. For \(n>1\)
there is one initial bulk transfer, one first right transfer and
\(2(n-1)\) loop transfers, totaling \(2n\). This count depends
on the specified movement granularity; it has not optimized arbitrary
program encodings or proved shortest description length.
\end{proof}
\end{annexitem}
